\VerdierPhysicalPageAnchor{0054}%
\VerdierSemanticAnchor{ch1-1-5-4-eq1-condition}
are $\alpha(\lambda,\mu)$-commutative when
$\lambda,\mu\in\mathbb Z^{n-1}$ or
$\lambda,\mu\in\mathbb Z e_n$. Let $Q(\sigma,\tau)$,
$\sigma\in\mathbb Z^{n-1}$, $\tau\in\mathbb Z e_n$, denote the
property:

\VerdierSemanticAnchor{ch1-1-5-4-property-q}
\begin{quote}
  \noindent\textup{$(Q(\sigma,\tau))$}\quad
  \emph{The following diagram is
  $\alpha(\sigma,\tau)\alpha^{-1}(\tau,\sigma)$-commutative:}
\end{quote}

\VerdierSemanticAnchor{ch1-1-5-4-eq2}
\begin{equation}
  \resizebox{0.98\textwidth}{!}{%
    $\begin{tikzcd}[ampersand replacement=\&,column sep=4.5em,row sep=3.2em]
      FT(\sigma+\tau)
        \arrow[r,"{F*c(\sigma,\tau)}"]
        \arrow[d,"{F*c(\tau,\sigma)}"]
        \& FT(\tau)T(\sigma)
        \arrow[r,"{p(\tau)*T(\sigma)}"]
        \& T'(\tau)FT(\sigma)
        \arrow[d,"{T'(\tau)*p(\sigma)}"'] \\
      FT(\sigma)T(\tau)
        \arrow[d,"{p(\sigma)*T(\tau)}"]
        \& {}
        \& T'(\tau)T'(\sigma)F
        \arrow[d,"{c'(\sigma,\tau)^{-1}*F}"'] \\
      T'(\sigma)FT(\tau)
        \arrow[r,"{T'(\sigma)*p(\tau)}"']
        \& T'(\sigma)T'(\tau)F
        \arrow[r,"{c'(\tau,\sigma)^{-1}*F}"']
        \& T'(\sigma+\tau)F
    \end{tikzcd}$%
  }\;.
  \tag*{\(\textup{(1.5.4.2)}_{\sigma,\tau}\)}
\end{equation}

By hypothesis, the properties
$Q(e_i,e_n)$, $Q(e_i,0)$, $1\leq i\leq n-1$, $Q(0,e_n)$
are true. Moreover, using the
$\alpha(\lambda,\mu)$-commutative diagrams
$\textup{(1.5.4.1)}_{\lambda,\mu}$, one readily verifies
(1.5.4.3):

\VerdierSemanticAnchor{ch1-1-5-4-eq3}
\begin{quote}
  \noindent\textup{(1.5.4.3)}\quad
  \emph{If two of the properties
  $Q(\sigma,\tau+\tau')$, $Q(\sigma,\tau)$, $Q(\sigma,\tau')$,
  $\sigma\in\mathbb Z^{n-1}$, $\tau,\tau'\in\mathbb Z e_n$
  hold, then so does the third. Likewise, if two of the properties
  $Q(\sigma+\sigma',\tau)$, $Q(\sigma,\tau)$,
  $Q(\sigma',\tau)$, $\sigma,\sigma'\in\mathbb Z^{n-1}$,
  $\tau\in\mathbb Z e_n$ hold, then so does the third.}
\end{quote}

\VerdierSemanticAnchor{ch1-1-5-4-q-conclusion}
It follows that the property $Q(\sigma,\tau)$ holds for every
$\sigma\in\mathbb Z^{n-1}$ and every $\tau\in\mathbb Z e_n$.

\VerdierSemanticAnchor{ch1-1-5-4-p-omega}
For every $\omega\in\mathbb Z^n$, write
$\omega=\sigma+\tau$, with $\sigma\in\mathbb Z^{n-1}$ and
$\tau\in\mathbb Z e_n$, and define the isomorphism
$p(\omega)\colon FT(\omega)\longrightarrow T'(\omega)F$ to be the
unique isomorphism making the diagram
$\textup{(1.5.4.1)}_{\sigma,\tau}$
$\alpha(\sigma,\tau)$-commutative. This enables us to write the
diagram $\textup{(1.5.4.1)}_{\lambda,\mu}$ for every pair
$\lambda,\mu\in\mathbb Z^n$ and to consider the property:

\VerdierSemanticAnchor{ch1-1-5-4-property-p}
\begin{quote}
  \noindent\textup{$(P(\lambda,\mu))$}\quad
  \emph{The diagram $\textup{(1.5.4.1)}_{\lambda,\mu}$ is
  $\alpha(\lambda,\mu)$-commutative.}
\end{quote}
